
% ---------- Feuille de route ----------
\begin{frame}{Suite du travail : feuille de route}
  \begin{block}{Trois chantiers pour la suite du projet}
    \begin{enumerate}
      \item \textbf{Implémentation des algorithmes} (recuit simulé, tabou,
      génétique) en Python.
      \item \textbf{Tests sur la base de données} réelle de la clinique.
      \item \textbf{Comparaison des performances} des trois métaheuristiques.
    \end{enumerate}
  \end{block}
  \vspace{0.2em}
  \begin{center}
    \begin{tikzpicture}[scale=0.9, every node/.style={font=\scriptsize}]
      \draw[thick, mcmPrimary, -{Latex}] (0,0) -- (11,0);
      \foreach \x/\lab in {0.5/Implémentation, 4.0/Tests, 7.5/Comparaison, 10.6/Livrable} {
        \draw[fill=mcmPrimary] (\x,0) circle (0.12);
        \node[above=2pt, align=center] at (\x,0) {\lab};
      }
    \end{tikzpicture}
  \end{center}
  \begin{alertblock}{\footnotesize Objectif}
    \footnotesize Passer des modèles mathématiques à un \textbf{prototype
    exécutable et mesuré}, utilisable comme base de décision.
  \end{alertblock}
\end{frame}

% ---------- Implémentation des algorithmes ----------
\begin{frame}{Implémentation des algorithmes}
  \begin{columns}[T]
    \begin{column}{0.55\textwidth}
      \small\textbf{Un moteur Python modulaire}
      \begin{itemize}
        \item Représentation commune d'un \textbf{planning} (solution)
        \item Fonction de coût reprenant l'objectif (lissage lits / ambu. / bloc)
        \item Un module par méthode : \texttt{recuit.py}, \texttt{tabou.py},
        \texttt{genetique.py}
        \item Même interface \texttt{solve(instance) -> planning} pour les
        comparer à conditions égales
      \end{itemize}
    \end{column}
    \begin{column}{0.45\textwidth}
      \centering
      \begin{tikzpicture}[scale=0.75]
        \draw[thick, mcmPrimary] (0,0) rectangle (3.4,0.6);
        \node[font=\scriptsize] at (1.7,0.3) {Instance (BDD)};
        \draw[-{Latex}] (1.7,0) -- (1.7,-0.7);
        \draw[thick, mcmExample] (0,-0.7) rectangle (3.4,-1.3);
        \node[font=\scriptsize] at (1.7,-1.0) {Fonction de coût};
        \draw[-{Latex}] (1.7,-1.3) -- (1.7,-2.0);
        \foreach \x/\m in {0/RS, 1.15/Tabou, 2.3/AG} {
          \draw[thick, mcmPrimaryMuted, fill=mcmPrimaryMuted] (\x,-2.0) rectangle (\x+1.0,-2.6);
          \node[font=\tiny] at (\x+0.5,-2.3) {\m};
        }
        \draw[-{Latex}] (1.7,-2.6) -- (1.7,-3.2);
        \draw[thick, mcmPrimary] (0,-3.2) rectangle (3.4,-3.8);
        \node[font=\scriptsize] at (1.7,-3.5) {Planning optimal};
      \end{tikzpicture}
    \end{column}
  \end{columns}
\end{frame}

% ---------- Tests sur la BDD ----------
\begin{frame}{Tests sur la base de données}
  \small
  \begin{columns}[T]
    \begin{column}{0.52\textwidth}
      \begin{block}{Données utilisées}
        \begin{itemize}
          \item Historique d'interventions de la clinique
          \item Durées opératoires et de séjour
          \item Part d'ambulatoire, vacations des chirurgiens
        \end{itemize}
      \end{block}
    \end{column}
    \begin{column}{0.48\textwidth}
      \begin{block}{Niveaux de test}
        \begin{itemize}
          \item \textbf{Unitaires} : contraintes C1--C5 respectées
          \item \textbf{Intégration} : lecture BDD $\rightarrow$ planning
          \item \textbf{Validation} : planning faisable et exploitable
        \end{itemize}
      \end{block}
    \end{column}
  \end{columns}
  \vspace{0.3em}
  \begin{exampleblock}{\footnotesize Jeux de tests}
    \footnotesize Une \textbf{instance jouet} (quelques patients) pour
    déboguer, puis des \textbf{instances réelles} de taille croissante pour
    éprouver le passage à l'échelle.
  \end{exampleblock}
\end{frame}


% ---------- Architecture technique ----------
\begin{frame}{Architecture technique : choix des frameworks}
  \begin{columns}[T]
    \begin{column}{0.48\textwidth}
      \begin{block}{Traitement -- Python / FastAPI}
        \small
        \begin{itemize}
          \item Moteur d'optimisation en \textbf{Python}
          \item \textbf{FastAPI} expose les algorithmes via une API REST
          \item Validation des données par \textbf{Pydantic}
          \item Calcul scientifique : \textbf{NumPy / Pandas}
        \end{itemize}
      \end{block}
    \end{column}
    \begin{column}{0.48\textwidth}
      \begin{block}{Interface web -- React}
        \small
        \begin{itemize}
          \item \textbf{React} (Vite) pour un tableau de bord interactif
          \item Visualisation des plannings et des performances
          \item Appels à l'API FastAPI
          \item Comparaison des méthodes en un coup d'œil
        \end{itemize}
      \end{block}
    \end{column}
  \end{columns}
  \vspace{0.2em}
  \begin{alertblock}{\footnotesize Pourquoi ce choix}
    \footnotesize FastAPI est rapide, typé et documenté automatiquement
    (Swagger) ; React offre une interface réactive et portable, séparant
    clairement le \textbf{moteur de calcul} de la \textbf{présentation}.
  \end{alertblock}
\end{frame}
